% ═══════════════════════════════════════════════════════════════════════════
% QR-CODE: Minitest Stromstärke & Spannung (Kl. 8, MT-01)
% Großer QR-Code auf einer A4-Seite (Beamer-Projektion)
% ═══════════════════════════════════════════════════════════════════════════

\documentclass[a4paper,11pt]{article}

\usepackage[utf8]{inputenc}
\usepackage[T1]{fontenc}
\usepackage{helvet}
\renewcommand{\familydefault}{\sfdefault}

\usepackage[margin=20mm]{geometry}
\usepackage{tikz}
\usepackage{xcolor}
\usepackage{qrcode}

% ═══════════════════════════════════════════════════════════════════════════
% KONFIGURATION
% ═══════════════════════════════════════════════════════════════════════════

\definecolor{physik}{HTML}{2c5aa0}
\colorlet{fachfarbe}{physik}

\newcommand{\sltitel}{Stromstärke \& Spannung}
\newcommand{\slurl}{https://devbox42.github.io/sciencesim/physik/kl08/elektrik/MT-01-stromstaerke-spannung.html}
\newcommand{\slurlkurz}{devbox42.github.io/.../MT-01-stromstaerke-spannung.html}

% ═══════════════════════════════════════════════════════════════════════════
% DOKUMENT
% ═══════════════════════════════════════════════════════════════════════════

\pagestyle{empty}

\begin{document}

\begin{center}

% Titel
\vspace*{10mm}
{\fontsize{28}{34}\selectfont\bfseries\color{fachfarbe}Minitest: \sltitel}

\vspace{15mm}

% QR-Code mit Rahmen
\fcolorbox{fachfarbe}{white}{\qrcode[height=100mm]{\slurl}}

\vspace{12mm}

% URL
{\large\color{gray}\slurlkurz}

\vspace{8mm}

% Hinweis
{\Large\color{gray}Scanne den QR-Code mit deinem Gerät}

\vspace{10mm}

% Info-Box
\fcolorbox{fachfarbe}{fachfarbe!10}{%
    \parbox{0.7\textwidth}{%
        \centering
        \vspace{3mm}
        {\large\bfseries Wichtig:} Wähle zuerst deine Platznummer (P1--P24)!\\[2mm]
        Arbeite konzentriert und lies die Aufgaben genau.
        \vspace{3mm}
    }%
}

\end{center}

\end{document}
